\graphicspath{ {./lab4/images/} }

\newpage
\section{Лабораторная работа № 4: Исследование стохастического возбуждения колебаний в ФХН-модели}
\subsection{Сравнение стохастических решений}

Проведем сравнение стохастических решений для $a = 1.1$, $a = 1.01$

(а) Построим графики $x(\varepsilon)$, пропустив переходный процесс $T = 100$ (Рис. \ref{fig:excitement_amplitude}).

\begin{figure}[htbp]
    \centering
    \includegraphics[width=1\textwidth]{excitement_amplitude.png}
    \caption{Амплитуда координаты $x$ при различных значениях $\varepsilon$.}
    \label{fig:excitement_amplitude}
\end{figure}

Локализуем зону стохастического возбуждения на основе 10 симуляций (Рис. \ref{fig:excitement_zone}).

\begin{figure}[htbp]
    \centering
    \includegraphics[width=1\textwidth]{excitement_zone.png}
    \caption{$\varepsilon$-зона стохастического возбуждения.}
    \label{fig:excitement_zone}
\end{figure}


(б) Построим временные ряды при различных значениях $\varepsilon$: с возбуждением и без возбуждения (Рис. \ref{fig:time_series_1_1}, \ref{fig:time_series_1_01}).

\begin{figure}[htbp]
    \centering
    \includegraphics[width=1\textwidth]{time_series_a=1,1.png}
    \caption{Временные ряды для $a=1.1$.}
    \label{fig:time_series_1_1}
\end{figure}

\begin{figure}[htbp]
    \centering
    \includegraphics[width=1\textwidth]{time_series_a=1,01.png}
    \caption{Временные ряды для $a=1.01$.}
    \label{fig:time_series_1_01}
\end{figure}



\newpage
\subsection{Доверительные эллипсы в допороговой и надпороговой зонах}

Построим доверительные эллипсы в допороговой и надпороговой зонах для $a = 1.1$, $a = 1.01$ (Рис. \ref{fig:confidence_ellipse_1_1}, \ref{fig:confidence_ellipse_1_01}).

В режиме без возбуждения доверительные эллипсы полностью ложатся в допороговую зону, и траектории совершают лишь небольшие колебания вокруг устойчивого равновесия.
В режиме возбуждения эллипсы захватывают надпороговую зону, и шум заставляет траектории выходить из окрестности равновесия на предельный цикл.
% Это особенно хорошо видно при $a = 1.01$.

\begin{figure}[htbp]
    \centering
    \includegraphics[width=1\textwidth]{confidence_ellipse_a=1,1_deterministic.png}
    \caption{Доверительные эллипсы для $a=1.1$.}
    \label{fig:confidence_ellipse_1_1}
\end{figure}

\begin{figure}[htbp]
    \centering
    \includegraphics[width=1\textwidth]{confidence_ellipse_a=1,01_deterministic.png}
    \caption{Доверительные эллипсы для $a=1.01$.}
    \label{fig:confidence_ellipse_1_01}
\end{figure}
